% !TeX root = ../main.tex

\chapter{相关理论概述}

\section{BYOL模型概述}

原版的BYOL模型处理对象为二维图像，目标是学习一个表示yθ，该表示随后可用于下游任务。BYOL使用两个神经网络进行学习：在线网络和目标网络。在线网络由一组权重 θ 定义，包含三个阶段：编码器（encoder）fθ，投影器（projector）gθ，预测器（predictor）qθ。目标网络具有与在线网络相同的架构，但使用另一组权重ξ。目标网络为训练在线网络提供回归目标（regression targets），其参数ξ是在线参数θ的指数移动平均值。更精确地说，给定一个目标衰减率（target decay rate）τ∈[0, 1]，在每次训练步骤后执行以下更新：ξ←τξ+(1-τ)θ。

输入与视图生成：给定一组图像D，一个从D中均匀采样的图像x，以及两个图像增强t和 t'。BYOL 通过分别应用图像增强t和 t'，从x生成两个增强视图：v ≜ t(x) 和 v' ≜ t'(x)。
前向传播与预测：从第一个增强视图v出发：在线网络输出一个表示yθ ≜ fθ(v)和一个投影 zθ ≜ gθ(y)。从第二个增强视图v'出发：目标网络输出一个表示 y'ξ ≜ fξ(v') 和一个目标投影 z'ξ ≜ gξ(y')。然后，在线网络输出目标投影 z'ξ 的一个预测 qθ(zθ)。预测器qθ仅应用于在线分支，这使得在线流程和目标流程之间的架构是不对称的。
损失函数：BYOL定义归一化后的预测与目标投影之间的均方误差为：
 
在每个训练步骤中执行随机优化步骤，仅针对参数θ最小化总损失.
参数更新：
BYOL 的动态更新过程总结如下：
 
训练结束时，仅保留编码器 fθ
本工作通过对BYOL模型进行修改，使其能够处理数据包长度序列数据，同时将编码器由resnet修改为GRU网络。GRU网络是循环神经网络的一种变体，在LSTM的基础上进行了简化，引入了更少的参数量和结构复杂度。GRU通过使用门控机制有效解决了传统RNN存在的梯度消失和梯度爆炸问题，适合处理长时间序列数据。BYOL模型中的两个变体分别为数据包长度序列原型及其变体。利用经过训练后的GRU网络对所有数据流的包长度序列进行映射，从而得到具有鲁棒性的流级别特征


\section{一维卷积神经网络概述}

IPV4网络中的网络层包头一般包含20字节的固定长度，其结构如上表所示。并且IP分组报头被进一步划分为以下12个字段，即，版本、首部长度、优先级与服务类型、总长度、标识、标志、段偏移、生存时间、协议、校验和、源IP地址和目的IP地址。源/目的IP地址与本地网络配置密切相关，但是与流量类型关系不大。从网络流量分类任务的角度来看，它们是冗余信息，会影响模型的泛化能力和分类结果。因此，本工作在提取数据过程中将其移除。
以往的方法通常直接将数据包数据作为输入特征，但是IP 数据包的字节序列结构比文本序列更固定且严谨。 连续的多个字节作为一个字段，为网络流量传输提供实际支持。例如，如上表所示，优先级与服务类型，总长度等都由多个字节组成。IP 数据包中，每个字节位置之间存在清晰的相对顺序。引入这种顺序信息有助于更好地表示数据包。因此，使用N-gram特征嵌入能够更好地学习字节序列中固有的结构和顺序信息，从而提取更加有效的包级别特征。考虑到数据处理的大小，本工作仅使用N = 2和N = 3的滑动窗口大小进行处理。
使用1-gram嵌入提取字节信息，2，3-gram嵌入提取数据包头内部的字节固定组合信息以及字节顺序信息。窗口大小分别为1、2、3，步长为1。处理过程如下图所示：



\section{多层感知机概述}

通过上述两个步骤，以及提取到了流级别特征和包级别特征，由于两个模态的特征对于分类的贡献度可能不同，因此需要使用特征融合策略联合学习二者的特征以实现精准鲁棒的流量分类。然而，现有融合策略（如直接拼接和加权拼接）多基于固定权重，没有考虑到不同特征之间的相互关系以及各个特征对分类的贡献程度，为此，本工作引入自适应深度特征融合（ADFF）机制。
对于其中一个特征（以流级别特征为例），ADFF首先计算生成三个向量：查询向量Qi，键向量Ki和值向量Vi，根据下述公式，可以计算流级别特征的查询向量分别与两个特征的键向量的点积，生成得分S1和S2，将得分通过softmax函数进行归一化生成权重，最后分别与两个特征的值向量相乘再求和即可得到对应的加权特征向量：(其中f表示流级别特征，p表示包级别特征)

在此之前已经做好了特征提取工作，因此多分类模型的选用多层感知机（MLP）。MLP通过ReLU等激活函数可学习复杂非线性决策边界，优于线性模型（如逻辑回归）和简单非线性模型（如SVM的线性核），能捕捉特征间的高阶交互关系，此外，相比于CNN/RNN等架构，没有不必要的卷积/循环计算开销，对高维特征向量处理效率高，因此可以更高效的进行分类任务。综上所述，本工作将使用MLP来进行加密流量的多分类任务。

\section{预处理}

本工作的技术路线可分为五步：第一步，收集网络流量数据集并将网络流量按照{源IP，目的IP，源端口号，目的端口号，使用协议}的五元组分割为数据流；第二步，从数据流中提取数据包长度序列，模拟TCP粘包拆包机制生成变体，使用BYOL模型提取变体和原型之间具有鲁棒性的流级别特征；第三步，从数据流中提取IP包头数据，使用n-gram技术和1D-CNN对其进行包级别特征提取；第四步，使用自适应深度特征融合技术(ADFF)将两个模态的特征进行融合，得到融合特征；第五步，将融合特征作为多分类模型的输入，进行多分类任务判断。其中，最关键的为二到四步：

流级别特征提取
通过模拟TCP的粘包拆包机制，构造数据包长度序列流量变体，并将同一流量的原型和变体进行适当的映射，使二者高维空间中距离拉近，从而提取出原型和变体之间具有鲁棒性的流级别特征。

包级别特征提取
数据包有效载荷会携带涉及用户隐私的信息，因此仅使用IP数据包头作为包级别特征的原始数据。通过n-gram技术和1D-CNN的训练提取出高效的包级别特征

特征融合
使用自适应深度特征融合技术(ADFF)，自适应地融合流级别特征和包级别特征，以解决两种特征对不同类型的流量分类贡献不同的问题，并通过各自的键向量和值向量，调整最终在高维向量所占的权重，从而将特征间的相互关系纳入计算

粘包：当应用程序通过 TCP连接 连续发送多个小的数据包时，TCP 为了提高网络效率可能会把这些 小的数据包合并成一个大的数据块进行发送，多个数据包之间没有没有明确的分隔，导致无法对这些数据包进行正确的读取，需要自行进行 数据边界识别，这种情况称为“粘包”。
拆包：反之，如果一次发送的数据量很大，TCP 可能会将其 划分为多个数据包进行发送，而在接收端这些数据包可能不会按照原来的边界一次性接收，而是 分散在多次接收操作中，否则会导致读取的数据不完整，这种情况称为“拆包”。

随机RTT生成：模拟真实网络中变化的往返时间；Nagle算法模拟：在RTT时间窗口内累积数据包，每个数据包随机生成一个到达延迟，消耗相应的网络延迟时间，模拟现实世界中的数据包到达间隔，累积结果存入缓冲区 buf；MSS分片处理：当缓冲区数据超过MSS时，分割成多个MSS大小的数据包，例如：buf=3000, MSS=1448 则输出1448和1448，剩余104；剩余数据处理：缓冲区剩余数据大于0但小于MSS时作为独立数据包输出，（类似TCP的Nagle算法行为）；序列标准化：对于处理后长度超过预设上限的数据流进行截断，对于处理后不足的数据流尾部进行零填充，使所有输出序列长度一致。

\section{论文章节安排}

\subsection{二级节标题}

\subsubsection{三级节标题}

本模板 \pkg{ustcthesis} 是中国科学技术大学本科生和研究生学位论文的 \LaTeX{}
模板, 按照《中国科学技术大学研究生学位论文撰写手册》（最近在修订中,以下简称《撰写手册》）和
《\href{https://www.teach.ustc.edu.cn/?attachment_id=13867}
{中国科学技术大学本科毕业论文（设计）格式}》的要求编写。

Lorem ipsum dolor sit amet, consectetur adipiscing elit, sed do eiusmod tempor
incididunt ut labore et dolore magna aliqua.
Ut enim ad minim veniam, quis nostrud exercitation ullamco laboris nisi ut
aliquip ex ea commodo consequat.
Duis aute irure dolor in reprehenderit in voluptate velit esse cillum dolore eu
fugiat nulla pariatur.
Excepteur sint occaecat cupidatat non proident, sunt in culpa qui officia
deserunt mollit anim id est laborum.


